\section{Discussion and Lessons Learned}

\subsection{Architectural Decisions}
Several architectural decisions shaped the final system:

\textbf{Three-layer separation with import-linter.} Enforcing strict layer boundaries at build time prevented coupling between Ingestion, Transformation, and Rendering. This paid dividends during the IR display model refactor, where the Transformation layer's data structures were restructured without touching any Rendering layer code. The \texttt{import-linter} tool caught accidental cross-layer imports on multiple occasions during development.

\textbf{Display model abstraction.} Rather than having Manim scenes query the scene graph directly, we introduced an explicit display model (the \texttt{TraceStep}/\texttt{RichTraceStep} and \texttt{AnimationCommand} types) that decouples rendering from data extraction. This made it straightforward to add the rich-SSA animation mode by extending the trace step type without modifying existing scene classes.

\textbf{Pipeline approach.} The tool processes data in a linear pipeline: parse $\rightarrow$ build scene graph $\rightarrow$ render or export. This simplified testing because each stage could be exercised independently with deterministic fixtures. The 302 tests in the suite are predominantly unit tests that operate on a single layer.

\textbf{llvmlite for terminator parsing.} Early prototypes used regex-based parsing of LLVM IR text, which proved fragile for edge cases such as \texttt{switch} and \texttt{invoke} terminators. Switching to \texttt{llvmlite}'s Python binding, which exposes the LLVM C~API, made CFG edge extraction reliable across all five terminator types.

\subsection{Deviations from Proposal}
The following proposed features were not implemented:

\begin{itemize}
    \item \textbf{Automatic C compilation:} The original proposal envisioned accepting \texttt{.c} files and invoking \texttt{clang} automatically. In practice, requiring users to supply pre-compiled \texttt{.ll} files simplified the tool's scope and avoided platform-specific compiler configuration issues.

    \item \textbf{Interactive stepping and pause:} Interactive playback was deprioritized in favor of richer static animation modes. Manim CE's rendering model produces video files rather than interactive sessions, making pause/step less natural.

    \item \textbf{YAML configuration and --palette:} Fine-grained customization via YAML files and named color palettes was deferred. The current CLI provides \texttt{--speed} and \texttt{--format} as the primary customization points.

    \item \textbf{GUI:} A graphical interface was listed as a stretch goal and was not pursued given the limited timeline.
\end{itemize}

Conversely, two capabilities were implemented that were \textit{not} in the original proposal:
\begin{itemize}
    \item \textbf{Rich-SSA three-column animation:} The \texttt{--ir-mode rich-ssa} mode displays symbolic SSA values alongside IR source and stack state, providing deeper instruction-level insight than originally~planned.
    \item \textbf{CFG traversal animation:} The \texttt{--cfg-animate} mode renders an animated walk through the control-flow graph using Graphviz DOT layout, a feature that emerged from explorations in the sandbox environment.
\end{itemize}

\subsection{Known Limitations}
The current implementation has several known limitations:

\begin{itemize}
    \item \textbf{Double-pop bug:} Under certain control-flow patterns involving consecutive returns, the basic animation mode may attempt to pop the stack twice for a single \texttt{ret} instruction. This is tracked as a known issue.

    \item \textbf{Symbolic-only SSA values:} The rich-SSA panel displays symbolic operand names (e.g., \texttt{\%3 = add \%1, \%2}) rather than concrete runtime values. A single swap-point in \texttt{ssa\_formatting.py} exists for when numeric runtime values are plumbed through in a future version.

    \item \textbf{No heap visualization:} The tool visualizes stack operations but does not yet render heap allocations or pointer-based data structures.

    \item \textbf{Manim render times:} Full-quality Manim renders can take significant time for programs with many instructions. The \texttt{--speed} flag and low-quality preview mode help during development, but render time remains a practical constraint for large programs.
\end{itemize}

\subsection{Lessons Learned}
\begin{itemize}
    \item \textbf{TDD was essential for refactors.} The test suite caught regressions immediately when refactoring the IR display model and restructuring the scene graph. Without comprehensive tests, these refactors would have been significantly riskier.

    \item \textbf{Architectural enforcement catches violations early.} \texttt{import-linter} rejected several pull requests that would have introduced cross-layer dependencies, preventing technical debt before it accumulated.

    \item \textbf{Sandbox-driven exploration.} Experimental Manim scripts in the~\texttt{sandbox/} directory allowed rapid prototyping of animation ideas (CFG traversal, register panels) without risking the main~codebase.

    \item \textbf{Typed CFG edges via llvmlite.} Moving from regex-based IR text parsing to \texttt{llvmlite}'s terminator API eliminated an entire class of parsing bugs and made the Ingestion layer robust against IR formatting variations.
\end{itemize}
